\documentclass[../../main.tex]{subfiles}
\begin{document}


{\bf [解]}
方程式
\begin{align}
x=\Bigl[\dfrac{1}{2}\Bigl(x+\dfrac{a}{x}\Bigr)\Bigr] \label{eq:1}
\end{align}
の右辺はガウス記号の性質および$x,a>0$より自然数だから，解$x$もまた自然数である．
そこで解を$x=k_a$（$k_a\in\mathbb N$）とおける．ガウス記号の性質から
\begin{align}
 & k_a\le\dfrac{1}{2}\Bigl(k_a+\dfrac{a}{k_a}\Bigr)<k_a+1 \\
 & k_a^2\le a<k_a^2+2k_a \\
\therefore
 & -1+\sqrt{a+1}<k_a\le\sqrt{a} \label{eq:2}
\end{align}
である．従って\cref{eq:2}を満たすような自然数$k_a$が存在すればそれが\cref{eq:1}の解である．

\subparagraph{(1)}

各$a=7,8,9$について個別に考える．

\subparagraph{$a=7$の時}

\cref{eq:2}で$a=7$として
\begin{align}
 -1+2\sqrt{2}<k_a\le\sqrt{7}
\end{align}
だから，$k_a=2$とすれば良い．

\subparagraph{$a=8$の時}

\cref{eq:2}で$a=8$として
\begin{align}
 2<k_a\le 2\sqrt{2}
\end{align}
だが，これを満たす自然数$k_a$はないため解なしである．

\subparagraph{$a=9$の時}

\cref{eq:2}で$a=9$として
\begin{align}
 -1+\sqrt{10}<k_a\le 3
\end{align}
だから，$k_a=3$とすれば良い．
\casesend

以上をまとめて，
\begin{align}
 \begin{cases}
  x = 2 & (a=7) \\
  \text{解なし} & (a=8) \\
  x = 3 & (a=9)
 \end{cases}
\end{align}
である．


\paragraph{(2)}

\cref{eq:2}について検討するため，$t\in\mathbb N$とする．
\begin{align}
 t^2 \le a < (t+1)^2 \\
\therefore
 t^2 \le a \le (t+1)^2-1 \quad (\because a \in\mathbb{N})
\end{align}
の時，
\begin{align}
 \left[\sqrt{a}\right] = t \label{eq:3}
\end{align}
だから，$a$が自然数であることに注意して\cref{eq:2}に代入すると
\begin{align}
 \sqrt{a+1}-1 < k_a \le t \label{eq:4}
\end{align}
である．そこで左辺について考える．

\subparagraph{$t^2 \le a < (t+1)^2-1$の時}

$a$が自然数だから
\begin{align}
 t^2 \le a \le (t+1)^2-2
\end{align}
である．この時
\begin{align}
 t^2+1 \le 1+a \le (t+1)^2-1 \\
\therefore
 & t < \sqrt{t^2+1} \le \sqrt{a+1} \le \sqrt{(t+1)^2-1} < \sqrt{(t+1)^2} = t+1 \\
\therefore
 & t-1 < \sqrt{a+1}-1 < t  
\end{align}
だから\cref{eq:4}を満たす$k_a$は$k_a=t$である．すなわち\cref{eq:2}が解を持つ．


\subparagraph{$a=(t+1)^2-1$の時}

この時は左辺が
\begin{align}
 \sqrt{a+1}-1 = \sqrt{(t+1)^2}-1 = t
\end{align}
となり，\cref{eq:4}を満たす自然数$k_a$が存在しない．すなわち\cref{eq:2}は解なしである．
\casesend

以上より\cref{eq:2}が解を持たないのは$a=(t+1)^2-1$（$t\in\mathbb N$）と表せるときであり，一般に
\begin{align}
 a_n=(n+1)^2-1 = n^2+2n \label{eq:5}
\end{align}
である．従って$n=1,2$として
\begin{align}
a_1=1^2+2\cdot1=3,\qquad a_2=2^2+2\cdot2=8
\end{align}
である．

\paragraph{(3)}

第$n$項までの和を
\begin{align}
 S_n=\sum_{k=1}^n\dfrac{1}{a_k} \label{eq:6}
\end{align}
とおくと，この$n\to\infty$の極限を求めれば良い．
\cref{eq:5}を代入すると
\begin{align}
 S_n
  &=\sum_{k=1}^n\dfrac{1}{k^2+2k} \\
  &=\sum_{k=1}^n\dfrac{1}{k(k+2)} \\
  &=\dfrac{1}{2}\sum_{k=1}^n\Bigl(\dfrac{1}{k}-\dfrac{1}{k+2}\Bigr) \\
  &=\dfrac{1}{2}\Bigl(1+\dfrac{1}{2}-\dfrac{1}{n+1}-\dfrac{1}{n+2}\Bigr) \\
  &\xrightarrow[n\to\infty]{} \dfrac{3}{4}
\end{align}
が求める極限値である．

\newpage

\end{document}
